% Part: first-order-logic
% Chapter: axiomatic-proofs
% Section: identity

\documentclass[../../../include/open-logic-section]{subfiles}

\begin{document}

\olfileid{fol}{axd}{ide}

\olsection{\usetoken{S}{identity}తో \usetoken{P}{derivation}}

\tetoken{వ్యుత్పత్తులలో}{!!{derivation}s} $\eq$ను చేర్చుకోవడానికి, మనం కొత్త స్వీకృత
పథకాలను మాత్రమే చేర్చుతాం. \tetoken{వ్యుత్పత్తి}{!!{derivation}}, $\Proves$ల నిర్వచనం
అదే ఉంటుంది; కొత్త స్వీకృతాలను కూడా అనుమతిస్తాం.

\begin{defn}[\tetoken{తాదాత్మ్య విధేయం}{!!{identity}} స్వీకృతాలు]
\begin{align}
\ollabel{ax:id1} & \eq[t][t], \\
\ollabel{ax:id2} & \eq[t_1][t_2] \lif (!B(t_1) \lif
  !B(t_2)),
\end{align}
ఇక్కడ $t$, $t_1$, $t_2$ ఏ మూస పదాలైనా కావచ్చు.
\end{defn}

\begin{prop}\ollabel{prop:sound}
  \olref{ax:id1}, \olref{ax:id2} స్వీకృతాలు సార్వత్రికంగా సత్యం.
\end{prop}

\begin{proof}
  సాధన.
\end{proof}

\begin{prob}
\olref[fol][axd][ide]{prop:sound}ను నిరూపించండి.
\end{prob}

\begin{prop}\ollabel{prop:iden1}
 ఏ మూస పదం~$t$, ఏ సమితి~$\Gamma$కైనా $\Gamma \Proves \eq[t][t]$.
\end{prop}
\sourcecorrection{OLTEAXD-011}{\texttt{ax:id1} మూస పదాలకే స్వీకృత
రూపాలను అనుమతిస్తుంది; కానీ ఆంగ్ల మూలంలోని ఈ ప్రతిపాదన ``ఏ పదం
$t$కైనా'' అని చెప్పింది. నిర్వచనానికి సరిపోయేలా ``ఏ మూస పదం
$t$కైనా'' అని సరిచేశాం.}

\begin{prop}\ollabel{prop:iden2}
  $\Gamma \Proves !A(t_1)$, $\Gamma \Proves \eq[t_1][t_2]$ అయితే,
  $\Gamma \Proves !A(t_2)$.
\end{prop}

\begin{proof}
\tetoken{సూత్రం}{!!{formula}}
\[
(\eq[t_1][t_2] \lif (!A(t_1) \lif !A(t_2)))
\]
అనేది~\olref{ax:id2} రూపం. ~\MPని రెండుసార్లు వర్తింపజేస్తే
నిష్కర్ష వస్తుంది.
\end{proof}

\end{document}
